\documentclass[10pt]{article}
\usepackage[T1,T2A]{fontenc}
\usepackage[utf8]{inputenc}
\usepackage[english,ukrainian]{babel}
\usepackage{geometry}
\usepackage{longtable}
\usepackage{array}
\usepackage{multirow}
\usepackage{hyperref}
\usepackage{mathptmx}
\renewcommand{\familydefault}{\rmdefault}

\geometry{a4paper,left=2cm,right=2cm,top=2cm,bottom=2cm}

\begin{document}

\begin{flushright}
ЗАТВЕРДЖЕНО\\
Наказом ТОВ "Криптографічні Телесистеми"\\
від \_\_ \_\_\_\_\_\_\_\_\_ 2026 р. № \_\_\_\\
(в редакції наказу ТОВ "Криптографічні Телесистеми"\\
від \_\_ \_\_\_\_\_\_\_\_\_ 2026 р. № \_\_)
\end{flushright}

\vspace{2cm}

\begin{center}
\textbf{\Large Телекомунікаційна система ERP/1}\\
\textbf{\Large «Повний набір сервісів»}\\
\vspace{1cm}
\textbf{\Large ЦІЛЬОВИЙ ПРОФІЛЬ БЕЗПЕКИ}\\
\vspace{0.5cm}
\large UA.46207985.СЗІ.ЦПБ-99
\end{center}

\newpage

\footnotesize
\begin{longtable}{|>{\centering\arraybackslash}p{0.5cm}|>{\raggedright\arraybackslash}p{4cm}|>{\raggedright\arraybackslash}p{5cm}|>{\raggedright\arraybackslash}p{2cm}|>{\raggedright\arraybackslash}p{3cm}|}
\hline
\textbf{№} & \textbf{Вимога з безпеки інформації} & \textbf{Вимога ГПБ} & \multicolumn{2}{c|}{\textbf{ЦПБ}} \\
\cline{4-5}
& & & \textbf{Захід захисту} & \textbf{Налаштований зміст заходу захисту}\\
\hline
\endfirsthead

\hline
\textbf{№} & \textbf{Вимога з безпеки інформації} & \textbf{Вимога ГПБ} & \multicolumn{2}{c|}{\textbf{ЦПБ}} \\
\cline{4-5}
& & & \textbf{Захід захисту} & \textbf{Налаштований зміст заходу захисту}\\
\hline
\endhead

\hline
\endfoot

\hline
\endlastfoot

\multicolumn{5}{|l|}{\textbf{1. УПРАВЛІННЯ ДОСТУПОМ (AC) (AC)}} \\ \hline
1 & Управління обліковими записами & Система забезпечує створення, активацію, зміну та видалення облікових записів відповідно до політик організації. Права доступу жорстко прив'язані до ідентифікатора працівника (`subject\_\allowbreak{}employee`) або ролі. & AC-2 & Параметр: account\_\allowbreak{}management\newline Тип: boolean\newline Значення: \textbf{true} \\ \hline
2 & Примусове застосування доступу (Access Enforcement) & Усі запити до ресурсів системи примусово проходять через Policy Decision Point (PDP) модуля `abac`. Прямий доступ до даних в обхід механізмів авторизації заборонено архітектурно. & AC-3 & Параметр: access\_\allowbreak{}enforcement\newline Тип: boolean\newline Значення: \textbf{true} \\ \hline
3 & Принцип найменших привілеїв (Least Privilege) & Профілі безпеки і ролі конфігуруються таким чином, що користувачам надаються лише ті привілеї, які необхідні для виконання їхніх посадових обов'язків. Механізм дозволів `abac` працює за принципом 'Default-Deny'. & AC-6 & Параметр: least\_\allowbreak{}privilege\newline Тип: boolean\newline Значення: \textbf{true} \\ \hline
4 & Блокування сеансу & Сеанс користувача автоматично блокується системою (`synrc/n2o` session expiry) після визначеного періоду бездіяльності, вимагаючи повторної автентифікації. & AC-11 & Параметр: session\_\allowbreak{}lock\_\allowbreak{}timeout\newline Тип: integer\newline Значення: \textbf{900} \\ \hline
5 & Завершення сеансу & Користувач може ініціювати завершення сеансу (logout). Крім того, система примусово завершує сеанси при досягненні максимального часу життя токена або зміни контексту безпеки. & AC-12 & Параметр: session\_\allowbreak{}termination\newline Тип: boolean\newline Значення: \textbf{true} \\ \hline
6 & Дозволені дії без ідентифікації чи автентифікації & Система явно визначає перелік відкритих сторінок та API-ендпойнтів (наприклад, портал входу або публічний довідник), доступних без автентифікації. Будь-які інші дії суворо заборонені. & AC-14 & Параметр: permitted\_\allowbreak{}unauthenticated\_\allowbreak{}actions\newline Тип: boolean\newline Значення: \textbf{true} \\ \hline
7 & Віддалений доступ & Віддалений доступ до системи керується додатковими політиками (наприклад, перевірка IP-адреси, VPN, mTLS), які інтегруються з `synrc/ca` та `abac` для забезпечення безпеки доступу за межами захищеного периметра. & AC-17 & Параметр: remote\_\allowbreak{}access\_\allowbreak{}control\newline Тип: boolean\newline Значення: \textbf{true} \\ \hline
8 & Динамічне пов'язання атрибутів (заборона кешування) & ABAC через компонент PIP (Policy Information Point) динамічно завантажує актуальні атрибути з контексту (наприклад, стан об'єкта `object\_\allowbreak{}process` чи профіль `subject\_\allowbreak{}employee` з KVS) безпосередньо в момент виконання запиту (`request`). Це гарантує, що рішення приймається на базі поточного, а не закешованого стану безпеки. & AC-16(1) & Параметр: dynamic\_\allowbreak{}binding\newline Тип: boolean\newline Значення: \textbf{true} \\ \hline
9 & Зміна значень атрибутів авторизованими особами через PAP & Точка адміністрування PAP дозволяє уповноваженим адміністраторам (які пройшли перевірку PDP) змінювати правила (`rule`), політики (`policy`) та кадрові атрибути співробітників, гарантуючи, що лише особи з належними повноваженнями можуть керувати системою доступу. & AC-16(2) & Параметр: authorized\_\allowbreak{}changes\newline Тип: boolean\newline Значення: \textbf{true} \\ \hline
10 & Підтримка системою пов'язання атрибутів (KVS) & Інфраструктура бази даних `synrc/kvs` нативно інтегрована з моделями `abac.hrl`. Система зберігає зв'язки між суб'єктами та атрибутами на рівні записів (records), що не підлягають несанкціонованому перезапису. & AC-16(3) & Параметр: attribute\_\allowbreak{}binding\_\allowbreak{}support\newline Тип: boolean\newline Значення: \textbf{true} \\ \hline
11 & Пов'язання атрибутів авторизованими особами & Тільки уповноважені особи (з відповідними атрибутами безпеки) можуть створювати або делегувати повноваження (наприклад, призначати ролі `assistant` або `delegate` в моделі `subject\_\allowbreak{}employee`). & AC-16(4) & Параметр: binding\_\allowbreak{}by\_\allowbreak{}authorized\_\allowbreak{}individuals\newline Тип: boolean\newline Значення: \textbf{true} \\ \hline
12 & Відображення атрибутів на пристроях виведення (nitro) & Завдяки інтеграції з `synrc/nitro`, рівні конфіденційності та безпекові атрибути об'єктів чітко рендеряться на інтерфейсі клієнта (веб-сторінках) у вигляді візуальних міток. & AC-16(5) & Параметр: display\_\allowbreak{}attributes\newline Тип: boolean\newline Значення: \textbf{true} \\ \hline
13 & Підтримка пов'язання атрибутів організацією (org, branch) & В `subject\_\allowbreak{}employee` підтримуються специфічні атрибути для ієрархічної структури (поля `org`, `branch`), що дозволяє обчислювати доступ на основі приналежності до організаційної одиниці. & AC-16(6) & Параметр: org\_\allowbreak{}binding\_\allowbreak{}support\newline Тип: boolean\newline Значення: \textbf{true} \\ \hline
14 & Послідовна інтерпретація атрибутів (XACML algorithms) & Використання стандартизованих алгоритмів комбінування політик XACML (алгоритми `all` або `any`) у PDP гарантує строгу та несуперечливу математичну інтерпретацію безпекових обмежень по всій системі. & AC-16(7) & Параметр: consistent\_\allowbreak{}interpretation\newline Тип: boolean\newline Значення: \textbf{true} \\ \hline
15 & Техніки та технології пов'язання атрибутів (Pattern Matching, CA) & Механізми Pattern Matching платформи Erlang/OTP використовуються для швидкого та безпомилкового порівняння атрибутів. Інтеграція криптографічних підписів (`synrc/ca`) у структуру об'єкта підтверджує автентичність його міток. & AC-16(8) & Параметр: binding\_\allowbreak{}techniques\newline Тип: boolean\newline Значення: \textbf{true} \\ \hline
16 & Перепризначення атрибутів згідно з життєвим циклом & Політики можуть передбачати логіку перепризначення атрибутів у процесі життєвого циклу об'єкта (наприклад, після підписання документа його атрибут статусу змінюється, що автоматично змінює правила доступу до нього). & AC-16(9) & Параметр: attribute\_\allowbreak{}reassignment\newline Тип: boolean\newline Значення: \textbf{true} \\ \hline
17 & Конфігурація атрибутів уповноваженими особами (PAP API) & Інфраструктура PAP надає API для гнучкого налаштування нових типів атрибутів і політик, гарантуючи, що ці налаштування виконуються тільки згідно встановлених процедур адміністрування. & AC-16(10) & Параметр: authorized\_\allowbreak{}configuration\newline Тип: boolean\newline Значення: \textbf{true} \\ \hline
18 & Атрибути безпеки об'єкту & BPE перед будь-яким переходом до наступного кроку (`bpe:next`) аналізує атрибути (метадані, документи), прив'язані до процесу. Шлюзи перевіряють ці мітки для ухвалення рішень маршрутизації. & AC-4(1) & Параметр: object\_\allowbreak{}security\_\allowbreak{}attributes\newline Тип: boolean\newline Значення: \textbf{true} \\ \hline
19 & Домени обробки (ізольовані процеси Erlang) & Кожен BPMN-процес виконується як окремий ізольований процес віртуальної машини Erlang (Actor). Ці процеси створюють незалежні безпекові домени обробки інформації без спільної пам'яті. & AC-4(2) & Параметр: processing\_\allowbreak{}domains\newline Тип: boolean\newline Значення: \textbf{true} \\ \hline
20 & Управління інформаційним потоком (SequenceFlow) & Сутність `sequenceFlow` (ребра графу BPMN) математично обмежує єдині можливі шляхи передачі інформації між задачами. & AC-4(3) & Параметр: enforce\_\allowbreak{}sequence\_\allowbreak{}flow\newline Тип: boolean\newline Значення: \textbf{true} \\ \hline
21 & Управління потоком зашифрованої інформації & BPE здатний маршрутизувати зашифровані payload-дані як непрозорі (opaque) структури, залишаючи розшифрування авторизованим клієнтам або конкретним сервісним завданням (`serviceTask`). & AC-4(4) & Параметр: encrypted\_\allowbreak{}payloads\newline Тип: boolean\newline Значення: \textbf{true} \\ \hline
22 & Вбудовування типів даних (Erlang Records) & Об'єкти інформаційного потоку строго типізовані через структури `record` (наприклад, `\#userTask\{\}`, `\#messageEvent\{\}`). Це запобігає переміщенню невідомих форматів. & AC-4(5) & Параметр: embedded\_\allowbreak{}data\_\allowbreak{}types\newline Тип: boolean\newline Значення: \textbf{true} \\ \hline
23 & Метадані (список docs у BPE) & BPE керує метаданими документів, використовуючи внутрішній список `docs` (API `bpe:docs/1`, `bpe:amend/2`), що невідривно супроводжує потік процесу. & AC-4(6) & Параметр: metadata\_\allowbreak{}management\newline Тип: boolean\newline Значення: \textbf{true} \\ \hline
24 & Механізми одностороннього потоку (DAG) & Направлені ациклічні графи (DAG) BPMN природно реалізують односторонню передачу даних від `beginEvent` до `endEvent` без можливості повернення (якщо це не передбачено схемою). & AC-4(7) & Параметр: one\_\allowbreak{}way\_\allowbreak{}flow\newline Тип: boolean\newline Значення: \textbf{true} \\ \hline
25 & Політики безпеки (Gateways) & Логічні шлюзи (`gateway`) виступають інтегрованими фільтрами, які застосовують політики безпеки на розгалуженнях потоку. & AC-4(8) & Параметр: security\_\allowbreak{}policies\newline Тип: boolean\newline Значення: \textbf{true} \\ \hline
26 & Перевірки, що проводить персонал (UserTask) & Використання `\#userTask\{\}` зупиняє автоматичний потік інформації та вимагає ручної перевірки (Human-in-the-loop) і підтвердження для продовження маршруту. & AC-4(9) & Параметр: personnel\_\allowbreak{}checks\newline Тип: boolean\newline Значення: \textbf{true} \\ \hline
27 & Активація та деактивація фільтрів політики (BoundaryEvent) & Використання подій, як-от `boundaryEvent` (Boundary Events), дозволяє динамічно активувати обхідні гілки політик в залежності від виникнення інцидентів або таймаутів. & AC-4(10) & Параметр: filter\_\allowbreak{}activation\newline Тип: boolean\newline Значення: \textbf{true} \\ \hline
28 & Конфігурація фільтрів політики безпеки (XML) & Правила маршрутизації (фільтри) завантажуються декларативно з BPMN XML файлів та транслюються у незмінні правила Erlang (`bpe:load/1`). & AC-4(11) & Параметр: filter\_\allowbreak{}configuration\newline Тип: boolean\newline Значення: \textbf{true} \\ \hline
29 & Ідентифікатори типу даних (MessageEvent) & Система розрізняє вхідні повідомлення через обробку `messageEvent` з відповідними ідентифікаторами повідомлень, ігноруючи нерозпізнані сигнали. & AC-4(12) & Параметр: data\_\allowbreak{}type\_\allowbreak{}identifiers\newline Тип: boolean\newline Значення: \textbf{true} \\ \hline
30 & Декомпозиція на субкомпоненти (Call Activity) & BPE підтримує виклик підпроцесів (`Call Activity`), що дозволяє розбити складні потоки на ізольовані субкомпоненти зі своїми власними політиками. & AC-4(13) & Параметр: subcomponent\_\allowbreak{}decomposition\newline Тип: boolean\newline Значення: \textbf{true} \\ \hline
31 & Обмеження фільтра політики (conditions) & Можливості маршрутизації обмежені набором умов `condition()`, які дозволяють математично гарантувати неперетин потоків. & AC-4(14) & Параметр: policy\_\allowbreak{}filter\_\allowbreak{}constraints\newline Тип: boolean\newline Значення: \textbf{true} \\ \hline
32 & Виявлення несанкціонованої інформації & При спробі передати процесу повідомлення, яке не відповідає поточній стадії потоку (unauthorized message), система автоматично відкидає його або фіксує в лог як невалідний event. & AC-4(15) & Параметр: unauthorized\_\allowbreak{}info\_\allowbreak{}detection\newline Тип: boolean\newline Значення: \textbf{true} \\ \hline
33 & Передача інформації про взаємопов'язані системи (ServiceTask) & Компонент `\#serviceTask\{\}` дозволяє інтегрувати механізми перевірки даних, перш ніж передати інформаційний потік до інших ІТС (наприклад, через API REST або RPC). & AC-4(16) & Параметр: interconnected\_\allowbreak{}systems\_\allowbreak{}transfer\newline Тип: boolean\newline Значення: \textbf{true} \\ \hline
34 & Прив'язка атрибуту безпеки (Process ID) & Інформація в рамках потоку жорстко прив'язується до Process ID (Id інстанції). Усі документи успадковують обмеження доступу поточного процесу. & AC-4(18) & Параметр: security\_\allowbreak{}attribute\_\allowbreak{}binding\newline Тип: boolean\newline Значення: \textbf{true} \\ \hline
35 & Перевірка метаданих (BPMN compare) & BPMN-умови підтримують функції типу `\{compare, Field, ConstCheckAgainst\}`, що перевіряють значення конкретних метаданих об'єкта під час прийняття рішення про маршрут. & AC-4(19) & Параметр: metadata\_\allowbreak{}validation\newline Тип: boolean\newline Значення: \textbf{true} \\ \hline
36 & Затверджені рішення (Default-Deny) & Erlang Match Specifications та XACML гарантують, що лише явно затверджені алгоритми маршрутизації виконуються (Default-Deny принцип). & AC-4(20) & Параметр: approved\_\allowbreak{}decisions\newline Тип: boolean\newline Значення: \textbf{true} \\ \hline
37 & Фізичне та логічне відділення інформаційних потоків & Логічне відділення гарантується механізмом процесів OTP VM. Можливість виконання різних процесів на різних фізичних нодах Erlang-кластера забезпечує фізичне відділення потоків. & AC-4(21) & Параметр: physical\_\allowbreak{}logical\_\allowbreak{}separation\newline Тип: boolean\newline Значення: \textbf{true} \\ \hline
38 & Єдиний доступ (bpe:next) & Централізований компонент контролю `bpe:next` слугує єдиною точкою (Choke Point) доступу та маршрутизації для всіх інформаційних потоків. & AC-4(22) & Параметр: single\_\allowbreak{}access\_\allowbreak{}point\newline Тип: boolean\newline Значення: \textbf{true} \\ \hline
39 & Модифікована інформація, яка не підлягає оприлюдненню (scrubbing) & При поверненні історії та статусів через API клієнту, BPE може відфільтровувати та очищати внутрішні атрибути (`scrubbing`), не допускаючи витоку чутливої інформації. & AC-4(23) & Параметр: scrub\_\allowbreak{}hidden\_\allowbreak{}fields\newline Тип: boolean\newline Значення: \textbf{true} \\ \hline
40 & Нормалізований формат (BPE Record) & Весь внутрішній потік конвертується в уніфікований стандартизований формат (`BPE Record Format`), незалежно від того, як дані надійшли на вході. & AC-4(24) & Параметр: normalized\_\allowbreak{}format\newline Тип: boolean\newline Значення: \textbf{true} \\ \hline
41 & Очищення даних (Garbage Collection, End Event) & Вбудовані механізми Garbage Collection віртуальної машини Erlang та знищення об'єктів пам'яті ізольованого процесу після завершення (`End Event`) забезпечують очищення даних з ОЗУ. & AC-4(25) & Параметр: data\_\allowbreak{}sanitization\newline Тип: boolean\newline Значення: \textbf{true} \\ \hline
42 & Дії з фільтрації аудиту (hist у KVS) & Кожен крок і зміна маршруту фіксується у вигляді об'єкта `hist` у `kvs`, створюючи прозорий ланцюг подій, де шлюзи та фільтри залишають свій слід. & AC-4(26) & Параметр: audit\_\allowbreak{}filter\_\allowbreak{}actions\newline Тип: boolean\newline Значення: \textbf{true} \\ \hline
43 & Незалежні фільтруючі механізми (ABAC + BPE) & Дозвіл на доступ (ABAC) та перевірка логіки процесу (BPE) працюють як архітектурно незалежні механізми, підсилюючи концепцію \textquotedbl{}Defense-in-depth\textquotedbl{}. & AC-4(27) & Параметр: independent\_\allowbreak{}filtering\newline Тип: boolean\newline Значення: \textbf{true} \\ \hline
44 & Лінійні фільтрувальні канали (Linear Pipeline) & BPE підтримує лінійну маршрутизацію (Linear Pipeline), коли `SequenceFlow` гарантовано не має відгалужень, забезпечуючи послідовне виконання політик без ризику обходу. & AC-4(28) & Параметр: linear\_\allowbreak{}pipelines\newline Тип: boolean\newline Значення: \textbf{true} \\ \hline
45 & Механізми оркестровки & Ядро BPE виступає єдиним централізованим механізмом оркестровки інформаційного обміну в рамках системи. & AC-4(29) & Параметр: orchestration\_\allowbreak{}mechanisms\newline Тип: boolean\newline Значення: \textbf{true} \\ \hline
46 & Фільтрація з використанням кількох процесів (Parallel Gateway) & Компонент `Parallel Gateway` (`\#gateway\{type=parallel\}`) створює декілька паралельних потоків, кожен з яких може мати власні незалежні правила фільтрації. & AC-4(30) & Параметр: multi\_\allowbreak{}process\_\allowbreak{}filtering\newline Тип: boolean\newline Значення: \textbf{true} \\ \hline
47 & Запобігання передачі вмісту & Якщо інформаційний об'єкт не відповідає політиці (фільтр на Gateway відхилив вміст), процес переривається або спрямовується на подію помилки (Error Event). & AC-4(31) & Параметр: content\_\allowbreak{}transfer\_\allowbreak{}prevention\newline Тип: boolean\newline Значення: \textbf{true} \\ \hline
48 & Вимоги до процесу передачі інформації (Receive/Send Task) & Механізми `Receive Task` та `Send Task` строго декларують вимоги до структури payload-повідомлень перед тим, як дозволити зовнішній чи внутрішній обмін. & AC-4(32) & Параметр: transfer\_\allowbreak{}process\_\allowbreak{}requirements\newline Тип: boolean\newline Значення: \textbf{true} \\ \hline
\multicolumn{5}{|l|}{\textbf{2. ІДЕНТИФІКАЦІЯ ТА АВТЕНТИФІКАЦІЯ (IA) (IA)}} \\ \hline
49 & Ідентифікація та автентифікація (організаційні користувачі) & Система однозначно ідентифікує та автентифікує (перевіряє справжність) організаційних користувачів та процеси, що діють від їх імені, перед тим, як дозволити їм доступ до захищених ресурсів. & IA-2 & Параметр: org\_\allowbreak{}user\_\allowbreak{}auth\newline Тип: boolean\newline Значення: \textbf{true} \\ \hline
50 & Ідентифікація та автентифікація (неорганізаційні користувачі) & Зовнішні користувачі, що підключаються з публічних мереж, проходять сувору автентифікацію (наприклад, за допомогою mTLS або інтеграції з КЕП) для запобігання анонімному або підробленому доступу. & IA-8 & Параметр: non\_\allowbreak{}org\_\allowbreak{}user\_\allowbreak{}auth\newline Тип: boolean\newline Значення: \textbf{true} \\ \hline
51 & Автентифікація на основі пароля (хешування) & Для паролів чи PING-кодів (наприклад, для розблокування контейнерів PKCS\#12) використовуються криптостійкі функції хешування з сіллю для запобігання атакам. & IA-5(1) & Параметр: password\_\allowbreak{}based\_\allowbreak{}auth\newline Тип: boolean\newline Значення: \textbf{true} \\ \hline
52 & Автентифікація на основі відкритого ключа & Є основним профілем бібліотеки `ca`. Вона генерує, перевіряє (validate) та керує життєвим циклом інфраструктури відкритих ключів (X.509), забезпечуючи строгу автентифікацію. & IA-5(2) & Параметр: pki\_\allowbreak{}enabled\newline Тип: boolean\newline Значення: \textbf{true} \\ \hline
53 & Зміна автентифікаторів до доставки (CSR) & Система підтримує протоколи, за яких приватний ключ генерується на пристрої клієнта, а до CA відправляється лише Certificate Signing Request (CSR). Це гарантує зміну/встановлення ключів ще до факту доставки готового сертифіката користувачу. & IA-5(5) & Параметр: pre\_\allowbreak{}delivery\_\allowbreak{}change\newline Тип: boolean\newline Значення: \textbf{true} \\ \hline
54 & Захист автентифікаторів (PEM/PKCS\#12) & Приватні ключі (автентифікатори) зберігаються в захищеному зашифрованому вигляді (наприклад, контейнери формату PEM з паролем або PKCS\#12), що унеможливлює використання ключів без знання пароля доступу. & IA-5(6) & Параметр: authenticator\_\allowbreak{}protection\newline Тип: boolean\newline Значення: \textbf{true} \\ \hline
55 & Відсутність вбудованих незашифрованих статичних автентифікаторів & Бібліотека заохочує динамічну генерацію ключів під час ініціалізації середовища (Root CA), що виключає необхідність зберігання \textquotedbl{}hardcoded\textquotedbl{} секретів у вихідному коді системи. & IA-5(7) & Параметр: no\_\allowbreak{}unencrypted\_\allowbreak{}static\newline Тип: boolean\newline Значення: \textbf{true} \\ \hline
56 & Багатосистемні облікові записи (mTLS/SSO) & Сертифікати X.509, видані `synrc/ca`, можуть використовуватись як багатосистемний засіб автентифікації (на кшталт mTLS/SSO) для доступу до різноманітних мікросервісів та ERP-компонентів, розгорнутих у спільному домені довіри. & IA-5(8) & Параметр: multi\_\allowbreak{}system\_\allowbreak{}accounts\newline Тип: boolean\newline Значення: \textbf{true} \\ \hline
57 & Управління об'єднанням автентифікаторів (Intermediate CA) & Підтримується архітектура федерації ключів (Federation) за рахунок використання ланцюгів сертифікатів (Intermediate CAs) та можливості перевіряти крос-сертифікаційні підписи (Cross-Certification). & IA-5(9) & Параметр: federation\_\allowbreak{}management\newline Тип: boolean\newline Значення: \textbf{true} \\ \hline
58 & Динамічне зв'язування мандатів & Бібліотека дозволяє видавати короткострокові сертифікати, що динамічно зв'язуються з конкретною сесією, IP-адресою чи атрибутами користувача (Dynamic Credential Binding). & IA-5(10) & Параметр: dynamic\_\allowbreak{}credential\_\allowbreak{}binding\newline Тип: boolean\newline Значення: \textbf{true} \\ \hline
59 & Автентифікація на основі апаратних токенів & `synrc/ca` прозоро підтримує роботу з підписами, згенерованими на апаратних захищених носіях (Hardware Tokens, НКІ, е-Токени, PKCS\#11). Серверній частині потрібен лише відкритий ключ для валідації. & IA-5(11) & Параметр: hardware\_\allowbreak{}token\_\allowbreak{}auth\newline Тип: boolean\newline Значення: \textbf{true} \\ \hline
60 & Ефективність біометричної автентифікації (Secure Enclave) & Хоча `ca` працює з криптографією, будь-яка біометрична автентифікація (наприклад, TouchID, FaceID) на клієнті конвертується в криптографічний підпис (наприклад, через Secure Enclave), який успішно валідується модулями `synrc/ca`. & IA-5(12) & Параметр: biometric\_\allowbreak{}efficiency\newline Тип: boolean\newline Значення: \textbf{true} \\ \hline
61 & Закінчення терміну дії автентифікаторів (NotBefore/NotAfter) & Жорстко перевіряються поля `NotBefore` та `NotAfter` в сертифікатах X.509. Сертифікати зі строком дії, що минув, автоматично відхиляються без можливості оскарження. & IA-5(13) & Параметр: authenticator\_\allowbreak{}expiration\newline Тип: boolean\newline Значення: \textbf{true} \\ \hline
62 & Управління змістом довірчих сховищ (Trust Stores) & На рівні сервера бібліотека керує сховищем довірених кореневих cертифікатів (Trust Store, `priv/cacerts`). Оновлення або відкликання (через CRL) довіри до Root CA автоматично реплікується на політику доступу. & IA-5(14) & Параметр: trust\_\allowbreak{}store\_\allowbreak{}management\newline Тип: boolean\newline Значення: \textbf{true} \\ \hline
63 & Продукти та послуги, затверджені уповноваженим органом (ДСТУ 4145-2002) & Завдяки абстракції криптографічних примітивів, `synrc/ca` підтримує державні криптографічні стандарти (наприклад, алгоритм цифрового підпису ДСТУ 4145-2002), необхідні для побудови систем, затверджених Державною службою спеціального зв'язку та захисту інформації (ДССЗЗІ). & IA-5(15) & Параметр: approved\_\allowbreak{}products\newline Тип: boolean\newline Значення: \textbf{true} \\ \hline
64 & Передача довірчої автентифікації зовнішньої сторони (АЦСК) & `synrc/ca` здатна перевіряти підписи, накладені зовнішніми Акредитованими центрами сертифікації ключів (АЦСК), забезпечуючи довірчу автентифікацію зовнішніх сторін. & IA-5(16) & Параметр: external\_\allowbreak{}auth\_\allowbreak{}transfer\newline Тип: boolean\newline Значення: \textbf{true} \\ \hline
65 & Менеджер паролів (KeyStores) & Незважаючи на те, що `synrc/ca` не є менеджером паролів у звичному розумінні, механізми захисту ключів (KeyStores, PKCS\#12) виступають у ролі захищених криптографічних сейфів, які можна відкрити лише спеціалізованими засобами із дотриманням секретності. & IA-5(18) & Параметр: password\_\allowbreak{}manager\newline Тип: boolean\newline Значення: \textbf{true} \\ \hline
\multicolumn{5}{|l|}{\textbf{3. АУДИТ ТА ПІДЗВІТНІСТЬ (AU) (AU)}} \\ \hline
66 & Узагальнення записів про аудит з декількох джерел (feeds) & BPMN підтримує паралельні ланцюги (feeds, streams), що дозволяє агрегувати записи про аудит з декількох джерел у єдиний потік журналу. & AU-2(1) & Параметр: audit\_\allowbreak{}aggregation\newline Тип: boolean\newline Значення: \textbf{true} \\ \hline
67 & Перегляд та оновлення (індексація hist) & Записи аудиту періодично індексуються, що дозволяє швидко витягувати всю історію (наприклад, через `bpe:hist/1`). & AU-2(3) & Параметр: audit\_\allowbreak{}review\_\allowbreak{}update\newline Тип: boolean\newline Значення: \textbf{true} \\ \hline
68 & Зміст записів аудиту (метадані, Permit/Deny) & Інформація в логах BPMN (особливо при збереженні об'єктів BPE/ABAC) містить вичерпні метадані: тип події, суб'єкт, цільовий об'єкт, результат операції (Permit/Deny) та позначку часу, що повністю закриває вимоги AU-3. & AU-3 & Параметр: audit\_\allowbreak{}content\_\allowbreak{}metadata\newline Тип: boolean\newline Значення: \textbf{true} \\ \hline
69 & Місткість сховища записів аудиту (розподілені бекенди) & BPMN може використовувати розподілені бекенди (Riak KV, розподілена Mnesia, KAI), що дозволяє горизонтально масштабувати місткість сховища логів (AU-4) без зупинки системи. & AU-4 & Параметр: audit\_\allowbreak{}storage\_\allowbreak{}capacity\newline Тип: boolean\newline Значення: \textbf{true} \\ \hline
70 & Позначка часу (erlang:system\_\allowbreak{}time) & Усі ідентифікатори та події, що генеруються в BPMN (наприклад, через генератор `id\_\allowbreak{}seq`), прив'язані до високоточного системного часу Erlang (`erlang:system\_\allowbreak{}time/1`). Це забезпечує надійну прив'язку подій до точного часу (AU-8(1)). & AU-8(1) & Параметр: sync\_\allowbreak{}time\_\allowbreak{}source\newline Тип: boolean\newline Значення: \textbf{true} \\ \hline
71 & Апаратні носії одноразового запису (Append-Only) & Архітектура Append-Only гарантує, що записи не оновлюються (no UPDATE). Це створює ефект апаратного \textquotedbl{}Write-Once-Read-Many\textquotedbl{} (AU-9(1)). Видалити запис без порушення криптографічної чи логічної цілісності ланцюга посилань практично неможливо. & AU-9(1) & Параметр: append\_\allowbreak{}only\_\allowbreak{}enforced\newline Тип: boolean\newline Значення: \textbf{true} \\ \hline
72 & Неспростовність (Ланцюжок збереження доказів) & Журнали BPMN слугують доказовою базою (Ланцюжок збереження доказів AU-10(3)). У комбінації з `synrc/ca` записи можуть підписуватися ЕЦП (AU-10(5)), що гарантує 100\% неспростовність дій автора транзакції. & AU-10(3) & Параметр: non\_\allowbreak{}repudiation\newline Тип: boolean\newline Значення: \textbf{true} \\ \hline
73 & Довгострокова можливість отримання записів аудиту & Сховища типу RocksDB (в якості бекенда BPMN) оптимізовані для довгострокового зберігання петабайтів даних (Long-Term Storage) із забезпеченням швидкого доступу до історичних записів (AU-11(1)). & AU-11(1) & Параметр: log\_\allowbreak{}retention\_\allowbreak{}days\newline Тип: integer\newline Значення: \textbf{1825} \\ \hline
74 & Огляд, аналіз та звітність про аудит & Аналіз аудиту відбувається за допомогою консольних інструментів платформи (наприклад, `kvs:hist/1` або інтерфейсів Erlang Shell). Також доступний експорт журналів для SIEM. & AU-6 & Параметр: audit\_\allowbreak{}analysis\_\allowbreak{}tools\newline Тип: boolean\newline Значення: \textbf{true} \\ \hline
75 & Генерація аудиту & KVS забезпечує генерацію записів аудиту для всіх операцій створення, оновлення (через створення нової ревізії) та видалення, гарантуючи фіксацію критичних подій безпеки. & AU-12 & Параметр: audit\_\allowbreak{}generation\newline Тип: boolean\newline Значення: \textbf{true} \\ \hline
\multicolumn{5}{|l|}{\textbf{4. ЦІЛІСНІСТЬ СИСТЕМИ ТА ІНФОРМАЦІЇ (SI) (SI)}} \\ \hline
76 & Перевірка вводу перевизначення (лише серверна валідація) & Клієнтська сторона (браузер) в `synrc/nitro` не містить логіки валідації, яку можна було б обійти чи перевизначити (override) через маніпуляції з DOM або JS. Уся перевірка жорстко зашита у виконуваний код Erlang (`synrc/form`). Жодне клієнтське перевизначення не здатне змінити правила серверної валідації. & SI-10(1) & Параметр: server\_\allowbreak{}side\_\allowbreak{}validation\_\allowbreak{}only\newline Тип: boolean\newline Значення: \textbf{true} \\ \hline
77 & Перевірка вводу інформації помилок (екранування HTML) & Повідомлення про помилки валідації генеруються на сервері (Erlang) і передаються клієнту через WebSocket як готові безпечні HTML-фрагменти (екрановані `nitro`), запобігаючи виникненню XSS під час відображення помилкового вводу користувачу. & SI-10(2) & Параметр: escape\_\allowbreak{}html\_\allowbreak{}output\newline Тип: boolean\newline Значення: \textbf{true} \\ \hline
78 & Передбачувана поведінка (декларативні Records) & Оскільки валідація описана декларативно (Record-структурами з типами полів, як-от `integer`, `string`, `date`), реакція системи на некоректні дані завжди є 100\% передбачуваною (Predictable Behavior). Відсутні приховані клієнтські скрипти, що можуть реагувати непередбачувано на специфічний ввід. & SI-10(3) & Параметр: predictable\_\allowbreak{}behavior\newline Тип: boolean\newline Значення: \textbf{true} \\ \hline
79 & Часові взаємодії (таймаути N2O) & Оскільки ввід передається через сталий WebSocket-зв'язок (N2O), система має вбудований контроль тайм-аутів з'єднання та захист від перевантажень (Rate-limiting). Довгі запити або неповний ввід просто ігноруються після завершення вікна сесії, блокуючи атаки типу Slowloris. & SI-10(4) & Параметр: timing\_\allowbreak{}interactions\newline Тип: boolean\newline Значення: \textbf{true} \\ \hline
80 & Обмеження вводу довіреними джерелами і форматами & Бібліотека `form` пропускає дані лише для тих полів, які явно задекларовані в схемі форми (Approved Formats). Pattern Matching платформи Erlang автоматично відкидає будь-які JSON/BERT payload-пакети, які містять незадекларовані поля (унеможливлюючи Mass Assignment Vulnerability) або походять з неавторизованого джерела (Trust Zone). & SI-10(5) & Параметр: strict\_\allowbreak{}type\_\allowbreak{}checking\newline Тип: boolean\newline Значення: \textbf{true} \\ \hline
81 & Профілактика вводу даних (блокування через ABAC) & Комбінація `nitro` та `abac` дозволяє здійснювати профілактику (Prevention): якщо користувач не має права вносити дані в певне поле, це поле взагалі не рендериться на стороні клієнта або рендериться в стані `readonly/disabled`, а бекенд-валідатор автоматично блокує будь-яку спробу його програмно передати (Data Input Prevention at source). & SI-10(6) & Параметр: data\_\allowbreak{}input\_\allowbreak{}prevention\newline Тип: boolean\newline Значення: \textbf{true} \\ \hline
\multicolumn{5}{|l|}{\textbf{5. ЗАХИСТ НОСІЇВ ІНФОРМАЦІЇ (MP) (MP)}} \\ \hline
82 & Доступ до носіїв інформації (через Erlang VM) & Прямий доступ до файлів БД (RocksDB/Mnesia) на носіях закритий на рівні ОС. Весь логічний доступ здійснюється виключно через інтерфейс KVS всередині віртуальної машини Erlang, яка діє як обмежений шлюз доступу (MP-2(1)). & MP-2(1) & Параметр: os\_\allowbreak{}level\_\allowbreak{}access\_\allowbreak{}control\newline Тип: boolean\newline Значення: \textbf{true} \\ \hline
83 & Криптографічний захист носіїв (Data-at-Rest) & KVS підтримує політики постійного збереження даних на диск (Disk Persistence). Захист інформації в стані спокою (Data-at-Rest) реалізується через інтеграцію з шифруванням файлової системи (LUKS) або шифруванням значень (MP-4(1)). & MP-4(1) & Параметр: encryption\_\allowbreak{}at\_\allowbreak{}rest\newline Тип: boolean\newline Значення: \textbf{true} \\ \hline
84 & Транспортування носіїв інформації (TLS-реплікація) & В розподіленому кластері KVS реплікує дані (переміщує їх) між нодами. Цей процес транспортування інформації захищається за допомогою TLS (Erlang Distribution over TLS), запобігаючи перехопленню даних поза контрольованими зонами (MP-5(1), MP-5(4)). & MP-5(1) & Параметр: tls\_\allowbreak{}replication\newline Тип: boolean\newline Значення: \textbf{true} \\ \hline
85 & Знищення інформації на носіях (Tombstones) & Хоча KVS є системою Append-Only, він надає чіткі API (`kvs:delete`) для виконання процедур очищення даних (наприклад, для забезпечення вимог GDPR \textquotedbl{}Право на забуття\textquotedbl{}). Під час видалення запису, бекенд RocksDB гарантовано помічає блоки як видалені (Tombstones), які під час подальшої компресії фізично стираються з носія без можливості відновлення (MP-6(8)). & MP-6(8) & Параметр: tombstone\_\allowbreak{}deletion\_\allowbreak{}enabled\newline Тип: boolean\newline Значення: \textbf{true} \\ \hline
86 & Використання носіїв інформації (марковані томи) & KVS може конфігуруватись на використання виключно заздалегідь визначених (іменованих) просторів таблиць на авторизованих томах зберігання, забороняючи використання невідомих чи немаркованих баз даних (MP-7(1)). & MP-7(1) & Параметр: authorized\_\allowbreak{}volumes\_\allowbreak{}only\newline Тип: boolean\newline Значення: \textbf{true} \\ \hline
\multicolumn{5}{|l|}{\textbf{6. ЗАХИСТ СИСТЕМ ТА КОМУНІКАЦІЙ (SC) (SC)}} \\ \hline
87 & Ізоляція функції безпеки & Функції безпеки (шифрування, авторизація) виконуються в ізольованих процесах Erlang VM, окремо від бізнес-логіки системи. Ця архітектура Actor Model унеможливлює втручання несанкціонованих процесів у роботу механізмів безпеки. & SC-2 & Параметр: security\_\allowbreak{}function\_\allowbreak{}isolation\newline Тип: boolean\newline Значення: \textbf{true} \\ \hline
88 & Інформація в загальних ресурсах системи & Erlang VM не використовує спільну пам'ять (shared memory) між процесами сесій N2O. Пам'ять звільняється автоматично (Garbage Collection) після обробки запиту, запобігаючи витоку залишкової інформації між різними користувачами. & SC-4 & Параметр: shared\_\allowbreak{}resource\_\allowbreak{}protection\newline Тип: boolean\newline Значення: \textbf{true} \\ \hline
89 & Захист від атак «відмова в обслуговуванні» (DoS) & Інфраструктура N2O/Cowboy обробляє сотні тисяч конкурентних WebSocket-з'єднань завдяки легкості Erlang-процесів. Вбудовані механізми rate-limiting, обмеження розміру payload та таймаути захищають від атак на виснаження ресурсів (Slowloris, Flood). & SC-5 & Параметр: dos\_\allowbreak{}protection\newline Тип: boolean\newline Значення: \textbf{true} \\ \hline
90 & Захист периметра (Boundary Protection) & Всі зовнішні підключення приймаються виключно через єдиний порт (керований інтерфейс), що обробляється веб-сервером (Cowboy). Доступ до внутрішніх API і вузлів кластера блокується правилами мережевого екранування за замовчуванням (Default Deny). & SC-7 & Параметр: boundary\_\allowbreak{}protection\newline Тип: boolean\newline Значення: \textbf{true} \\ \hline
91 & Захист цілісності та конфіденційності під час передачі & Передача інформації між клієнтом і сервером здійснюється виключно через зашифровані канали (WSS/TLS), що унеможливлює перехоплення (Man-in-the-Middle) або несанкціоновану модифікацію даних на льоту. & SC-8 & Параметр: transmission\_\allowbreak{}confidentiality\newline Тип: boolean\newline Значення: \textbf{true} \\ \hline
92 & Розрив мережевого з'єднання & Протокол N2O автоматично розриває мережеве з'єднання після завершення комунікаційного сеансу або перевищення часу очікування (Inactivity Timeout), звільняючи ресурси та сесію. & SC-10 & Параметр: network\_\allowbreak{}disconnect\_\allowbreak{}timeout\newline Тип: integer\newline Значення: \textbf{3600} \\ \hline
93 & Криптографічне управління ключами & Керування ключами та сертифікатами для встановлення захищених з'єднань делегується бібліотеці `synrc/ca` та інтегрованим механізмам TLS ОС, що відповідають чинним криптографічним стандартам. & SC-12 & Параметр: cryptographic\_\allowbreak{}key\_\allowbreak{}management\newline Тип: boolean\newline Значення: \textbf{true} \\ \hline
94 & Криптографічний захист & Система підтримує використання схвалених криптографічних алгоритмів (AES-GCM, SHA-256/512, ДСТУ 4145-2002 через відповідні модулі), які забезпечують необхідний рівень стійкості. & SC-13 & Параметр: cryptographic\_\allowbreak{}protection\newline Тип: boolean\newline Значення: \textbf{true} \\ \hline
95 & Автентичність сеансу зв'язку & Автентичність WebSocket-сесій N2O захищається токенами з криптографічним підписом (HMAC). Це гарантує, що повідомлення в рамках сесії не можуть бути підроблені або інжектовані сторонньою особою. & SC-23 & Параметр: session\_\allowbreak{}authenticity\newline Тип: boolean\newline Значення: \textbf{true} \\ \hline
\multicolumn{5}{|l|}{\textbf{7. УПРАВЛІННЯ КОНФІГУРАЦІЄЮ (CM) (CM)}} \\ \hline
96 & Базова конфігурація (Baseline Configuration) & Файл `erpuno.ex` фіксує єдину базову конфігурацію (Baseline) безпеки для всієї системи. Зміни в конфігурації застосовуються виключно через оновлення коду, що забезпечує строгий контроль та аудит змін. & CM-2 & Параметр: baseline\_\allowbreak{}configuration\_\allowbreak{}code\newline Тип: boolean\newline Значення: \textbf{true} \\ \hline
97 & Контроль змін конфігурації & Оскільки CMDB є частиною вихідного коду (Elixir модуль), всі зміни проходять через стандартний процес розробки (Code Review, CI/CD, контроль версій Git), унеможливлюючи несанкціоноване втручання в налаштування на \textquotedbl{}живій\textquotedbl{} системі. & CM-3 & Параметр: configuration\_\allowbreak{}change\_\allowbreak{}control\newline Тип: boolean\newline Значення: \textbf{true} \\ \hline
98 & Налаштування конфігурації & Параметри політик безпеки та їхні значення за замовчуванням жорстко зашиті в `erpuno.ex`. Компілятор генерує незмінний об'єктний файл, який гарантує виконання затверджених налаштувань під час роботи. & CM-6 & Параметр: compile\_\allowbreak{}time\_\allowbreak{}config\newline Тип: boolean\newline Значення: \textbf{true} \\ \hline
99 & Інвентаризація компонентів ІТС & Профіль автоматично фіксує перелік використовуваних компонентів системи (бібліотеки `synrc/kvs`, `synrc/ca` тощо) та їхні параметри, забезпечуючи інтегральну інвентаризацію модулів захисту. & CM-8 & Параметр: component\_\allowbreak{}inventory\newline Тип: boolean\newline Значення: \textbf{true} \\ \hline
100 & План управління конфігурацією & Профіль `ERPUNO.Profile` слугує самодокументованим планом управління конфігураціями, здатним автоматично генерувати актуальну документацію (через метод `generate\_\allowbreak{}md/0`). & CM-9 & Параметр: auto\_\allowbreak{}generate\_\allowbreak{}documentation\newline Тип: boolean\newline Значення: \textbf{true} \\ \hline
\multicolumn{5}{|l|}{\textbf{8. ПЛАНУВАННЯ БЕЗПЕРЕРВНОЇ РОБОТИ (CP) (CP)}} \\ \hline
101 & Резервне копіювання інформаційної системи & Система здійснює автоматичне створення резервних копій бази даних KVS. Бекапи можуть зберігатися локально або відправлятися до хмарних сховищ, гарантуючи відновлення даних відповідно до RTO/RPO. & CP-9 & Параметр: system\_\allowbreak{}backup\newline Тип: boolean\newline Значення: \textbf{true} \\ \hline
102 & Відновлення інформаційної системи & KVS забезпечує механізми швидкого відновлення з резервних копій (`kvs:restore/1`), дозволяючи адміністраторам відновити працездатність вузла з гарантією логічної цілісності журналу транзакцій. & CP-10 & Параметр: system\_\allowbreak{}recovery\newline Тип: boolean\newline Значення: \textbf{true} \\ \hline
\multicolumn{5}{|l|}{\textbf{9. РЕАГУВАННЯ НА ІНЦИДЕНТИ (IR) (IR)}} \\ \hline
103 & Політика та процедури реагування на інциденти & Політика реагування на інциденти закодована безпосередньо у бізнес-процесах (BPMN) `erpuno/itsm`, забезпечуючи неухильне виконання процедур на кожному кроці життєвого циклу інциденту. & IR-1 & Параметр: incident\_\allowbreak{}policy\_\allowbreak{}enforcement\newline Тип: boolean\newline Значення: \textbf{true} \\ \hline
104 & Навчання з реагування на інциденти & Інтерфейс `erpuno/itsm` надає контекстні підказки та інструкції безпосередньо у формах обробки інцидентів, що виконує функцію інтерактивного навчання та підтримки персоналу в кризових ситуаціях. & IR-2 & Параметр: incident\_\allowbreak{}response\_\allowbreak{}training\newline Тип: boolean\newline Значення: \textbf{true} \\ \hline
105 & Тестування реагування на інциденти & Можливості BPE-движка дозволяють запускати симуляції інцидентів у тестовому середовищі `itsm` для регулярного тестування планів реагування без впливу на продуктивну систему. & IR-3 & Параметр: incident\_\allowbreak{}testing\_\allowbreak{}simulation\newline Тип: boolean\newline Значення: \textbf{true} \\ \hline
106 & Обробка інцидентів & Будь-який інцидент безпеки автоматично реєструється як тикет у модулі `itsm`. Процес обробки (`incident\_\allowbreak{}management\_\allowbreak{}process`) гарантує, що інцидент пройде всі необхідні етапи: виявлення, ізоляція, усунення, відновлення. & IR-4 & Параметр: incident\_\allowbreak{}handling\newline Тип: boolean\newline Значення: \textbf{true} \\ \hline
107 & Моніторинг інцидентів & Статуси інцидентів у реальному часі відображаються на дашбордах `erpuno/itsm`. BPE-движок дозволяє призначати SLA-таймаути (Boundary Events) для ескалації інцидентів, якщо час реакції перевищено. & IR-5 & Параметр: incident\_\allowbreak{}monitoring\newline Тип: boolean\newline Значення: \textbf{true} \\ \hline
108 & Звітність про інциденти & Система `erpuno/itsm` автоматично формує звіти про інциденти, надаючи вичерпну інформацію про порушені сервіси, час простою (SLA) та дії команди, що закриває вимоги внутрішньої та зовнішньої звітності. & IR-6 & Параметр: automated\_\allowbreak{}incident\_\allowbreak{}reporting\newline Тип: boolean\newline Значення: \textbf{true} \\ \hline
109 & Допомога у реагуванні на інциденти & Вбудовані засоби комунікації в рамках тикету (коментарі, призначення експертів, інтеграція з базою знань) забезпечують оперативну допомогу та координацію групи безпеки під час інциденту. & IR-7 & Параметр: incident\_\allowbreak{}assistance\newline Тип: boolean\newline Значення: \textbf{true} \\ \hline
110 & План реагування на інциденти & Сам конфігураційний код процесів `erpuno/itsm` виступає в ролі актуального та виконуваного плану реагування на інциденти (Executable Incident Response Plan), усуваючи розбіжність між документацією та реальністю. & IR-8 & Параметр: executable\_\allowbreak{}ir\_\allowbreak{}plan\newline Тип: boolean\newline Значення: \textbf{true} \\ \hline
111 & Реагування на витік інформації (Information Spillage) & Спеціалізовані BPMN-схеми в `itsm` передбачені для обробки витоків інформації. Вони включають обов'язкові кроки ідентифікації скомпрометованих даних, блокування доступу (`abac`) та сповіщення регулятора. & IR-9 & Параметр: information\_\allowbreak{}spillage\_\allowbreak{}response\newline Тип: boolean\newline Значення: \textbf{true} \\ \hline
\multicolumn{5}{|l|}{\textbf{10. ОЦІНКА РИЗИКІВ (RA) (RA)}} \\ \hline
112 & Оцінка ризиків & Всі ідентифіковані вразливості або загрози фіксуються в реєстрі `itsm`, де їм присвоюється рівень критичності. На базі цього рівня розраховується пріоритет та призначаються відповідальні особи. & RA-3 & Параметр: risk\_\allowbreak{}assessment\newline Тип: boolean\newline Значення: \textbf{true} \\ \hline
113 & Сканування вразливостей & Система передбачає можливість інтеграції зовнішніх сканерів вразливостей, звіти яких автоматично перетворюються на тикети усунення ризиків в системі `itsm`. & RA-5 & Параметр: vulnerability\_\allowbreak{}scanning\newline Тип: boolean\newline Значення: \textbf{true} \\ \hline
\multicolumn{5}{|l|}{\textbf{11. ФІЗИЧНИЙ ЗАХИСТ ТА ЗАХИСТ НАВКОЛИШНЬОГО СЕРЕДОВИЩА (PE) (PE)}} \\ \hline
114 & Фізичний доступ & Регламентується інструкціями з режиму ЦОД (див. розділ PE у digital-profile.pdf). & PE-2 & Параметр: physical\_\allowbreak{}access\_\allowbreak{}doc\_\allowbreak{}ref\newline Тип: string\newline Значення: \textbf{digital-profile.pdf} \\ \hline
115 & Контроль фізичного доступу & Регламентується системами СКУД будівлі (див. розділ PE у digital-profile.pdf). & PE-3 & Параметр: physical\_\allowbreak{}access\_\allowbreak{}control\_\allowbreak{}ref\newline Тип: string\newline Значення: \textbf{digital-profile.pdf} \\ \hline
\multicolumn{5}{|l|}{\textbf{12. ОБІЗНАНІСТЬ ТА НАВЧАННЯ (AT) (AT)}} \\ \hline
116 & Навчання з обізнаності щодо безпеки & Програми навчання та інструктажі (див. розділ AT у digital-profile.pdf). & AT-2 & Параметр: security\_\allowbreak{}awareness\_\allowbreak{}training\_\allowbreak{}ref\newline Тип: string\newline Значення: \textbf{digital-profile.pdf} \\ \hline
117 & Навчання з питань безпеки, залежне від ролі & Спеціалізоване навчання для адміністраторів (див. розділ AT у digital-profile.pdf). & AT-3 & Параметр: role\_\allowbreak{}based\_\allowbreak{}training\_\allowbreak{}ref\newline Тип: string\newline Значення: \textbf{digital-profile.pdf} \\ \hline
\multicolumn{5}{|l|}{\textbf{13. УПРАВЛІННЯ РИЗИКАМИ В ЛАНЦЮГУ ПОСТАЧАННЯ (SR) (SR)}} \\ \hline
118 & План управління ризиками ланцюга постачання & План управління базується на принципі відмови від стороннього коду. Замість інтеграції сотень неперевірених бібліотек, команда розробляє всі необхідні компоненти (веб-сервер N2O, KVS, CA) самостійно в єдиному монорепозиторії. & SR-2 & Параметр: zero\_\allowbreak{}dependency\_\allowbreak{}policy\newline Тип: boolean\newline Значення: \textbf{true} \\ \hline
119 & Контроль та захист ланцюга постачання & Заборонено використання автоматичних менеджерів пакетів для завантаження коду під час збірки. Весь вихідний код системи та модулів міститься безпосередньо в репозиторії проекту, що гарантує 100\% контроль над кожним рядком коду, який компілюється. & SR-3 & Параметр: strict\_\allowbreak{}code\_\allowbreak{}vendoring\newline Тип: boolean\newline Значення: \textbf{true} \\ \hline
120 & Аудит та огляди постачальників & Єдиним зовнішнім «постачальником» є розробник платформи Erlang/OTP (компанія Ericsson та спільнота). Оновлення версій OTP проходять жорстку перевірку перед їх застосуванням на продуктивних серверах. & SR-6 & Параметр: erlang\_\allowbreak{}otp\_\allowbreak{}supplier\_\allowbreak{}review\newline Тип: boolean\newline Значення: \textbf{true} \\ \hline
121 & Автентичність компонентів & Будь-які оновлення або патчі самої системи передаються виключно через захищені канали із застосуванням криптографічних підписів (`synrc/ca`), що виключає можливість підміни бінарних файлів під час розгортання. & SR-11 & Параметр: component\_\allowbreak{}authenticity\newline Тип: boolean\newline Значення: \textbf{true} \\ \hline
\multicolumn{5}{|l|}{\textbf{14. ПЛАНУВАННЯ (PL) (PL)}} \\ \hline
122 & План безпеки системи & Система підтримує автоматичну генерацію актуального Плану безпеки системи на основі цього профілю (CM). План чітко визначає всі технічні та організаційні контролі (AC, AU, SC, CM тощо), середовище функціонування та архітектуру. & PL-2 & Параметр: system\_\allowbreak{}security\_\allowbreak{}plan\newline Тип: boolean\newline Значення: \textbf{true} \\ \hline
123 & Правила поведінки & Для всіх користувачів платформи (організаційних та зовнішніх) визначаються правила поведінки. Згода з правилами (наприклад, Terms of Service, NDA) фіксується в базі даних KVS під час онбордингу, перед наданням доступу до ERP/1. & PL-4 & Параметр: rules\_\allowbreak{}of\_\allowbreak{}behavior\newline Тип: boolean\newline Значення: \textbf{true} \\ \hline
124 & Архітектура безпеки & Архітектура безпеки платформи (Actor Model, ізоляція процесів, Zero Dependencies, Append-Only логування) розробляється інтегровано з загальною архітектурою підприємства. Вона регулярно переглядається на відповідність державним стандартам (НД ТЗІ, NIST SP 800-53). & PL-8 & Параметр: security\_\allowbreak{}architecture\newline Тип: boolean\newline Значення: \textbf{true} \\ \hline


\end{longtable}

\end{document}
